% Generated draft: requires mathematical and source review.
\documentclass[11pt]{article}
\usepackage[margin=25mm]{geometry}
\usepackage{fontspec}
\setmainfont{Latin Modern Roman}
\usepackage{amsmath,amssymb,mathtools,unicode-math}
\setmathfont{Latin Modern Math}
\usepackage{enumitem,needspace,xurl,hyperref}
\setlength{\parindent}{0pt}
\setlength{\parskip}{5pt}
\setlength{\emergencystretch}{3em}
\newcommand{\N}{\mathbb{N}}
\newcommand{\Z}{\mathbb{Z}}
\newcommand{\Q}{\mathbb{Q}}
\newcommand{\R}{\mathbb{R}}
\newcommand{\C}{\mathbb{C}}
\newcommand{\F}{\mathbb{F}}
\newcommand{\A}{\mathbb{A}}
\newcommand{\PP}{\mathbb{P}}
\newcommand{\E}{\mathbb{E}}
\newcommand{\eps}{\varepsilon}
\newcommand{\dd}{\,\mathrm{d}}
\begin{document}
\section*{UnsolvedMath ID: 701001}
\textbf{Source problem code:} catalogue entry
\subsection{Definitions and mathematical statement}
Let \(G\) be a group and let \(V\) be any set of group words, each involving finitely many variables. Define the set of word values and their inverses by
\[
W_V(G)=\{w(g_1,\ldots,g_k)^{\varepsilon}: w\in V,\ w\text{ involves }k\text{ variables},\ g_1,\ldots,g_k\in G,\ \varepsilon\in\{1,-1\}\}.
\]
The verbal subgroup associated with \(V\) is \(V(G)=\langle W_V(G)\rangle\). Following the attribution to Yu. I. Merzlyakov in the supplied source, its width with respect to \(V\) is finite if there exists a nonnegative integer \(m\) such that every \(h\in V(G)\) admits an expression
\[
h=x_1\cdots x_r,\qquad 0\le r\le m,\qquad x_i\in W_V(G).
\]
The empty product is the identity. The width is the least such bound, subject to the convention for whether zero is included among natural numbers, and is \(\infty\) if no finite bound exists. The value \(\infty\) indicates unbounded finite product lengths, not infinite products.

For a finitely generated group \(G\), choose a finite symmetric generating set \(S\), and let
\[
\gamma_{G,S}(n)=\bigl|\{g\in G: |g|_S\le n\}\bigr|,\qquad n\in\mathbb Z_{\ge0},
\]
where \(|g|_S\) is the least number of generators from \(S\) whose product is \(g\). The group has polynomial growth if \(\gamma_{G,S}(n)\le C(n+1)^d\) for all \(n\), for some \(C>0\) and nonnegative integer \(d\). It has intermediate growth if its growth is subexponential,
\[
\lim_{n\to\infty}\gamma_{G,S}(n)^{1/n}=1,
\]
but it has no polynomial growth bound. These growth classifications do not depend on the chosen finite generating set.

\textbf{Problem (V. G. Bardakov, 1999; catalogue entry).} Is it true that for every finitely generated group \(G\) of intermediate growth and every set \(V\) of group words, the verbal subgroup \(V(G)\) has finite width with respect to \(V\)? The bound may depend on both \(G\) and \(V\).

The supplied source states as background that the analogous finite-width assertion holds for finitely generated groups of polynomial growth. This background assertion and the supplied open status are retained as source claims; neither has been independently verified here.
\subsection{Short English statement}
Does every set of group words have a finite product-length bound in every finitely generated group of intermediate growth? More precisely, must every element of the associated verbal subgroup be a product of a uniformly bounded number of word values or their inverses, with the bound allowed to depend on the group and the set of words?
\subsection{Sources}
\begin{itemize}
\item Local UnsolvedMath registry, supplied complete source record: permanent ID \(701001\), exact source problem code \texttt{catalogue entry}, title \emph{Kourovka Notebook Problem 14.16}. The registry labels the problem open; this is a supplied status, not an independently verified assessment.
\item \emph{The Kourovka Notebook (Unsolved Problems in Group Theory)}, 14th issue (1999), Problem 14.16; proposed by V. G. Bardakov. The supplied record also identifies the 21st issue edition (2026), \texttt{arXiv:1401.0300v48}. The collection description names E. I. Khukhro and V. D. Mazurov as editors.
\item The supplied statement attributes the definition of verbal-subgroup width to Yu. I. Merzlyakov and reports the polynomial-growth finiteness assertion as known. No separate reference for either attribution or assertion is supplied.
\end{itemize}
\end{document}
